\section{Example: Parallel Processing} \label{ex10}

There are two ways to parallel processing in WAVE. The first is for the
case of a single electron, i.e.

\cboxs{IBUNCH=0}

This mode makes use of the
OpenMP features of the FORTRAN compiler. To activate is set

\cboxs{MTHREADS= number of CPU cores to use}

For

\cboxs{MTHREADS $<$ 0 }

all cores are used. You can explore this option e.g. by changing the
setting of the \ref{ex4} example accordingly.

The second way for the multi-particle mode, i.e.

\cboxs{IBUNCH=1}

or

\cboxs{IBUNCH=-1}

is to spawn several instances of WAVE by setting

\cboxs{ICLUSTER=1}

If this mode is active the maximum of \scap{ICLUSTER} and
\scap{MTHREADS} gives the number of the spawned instances. Try
it by modifying the \ref{ex9}~ example.\\

